\section{Sphere maps, framings and homotopy classes}\label{sec:III:maps}
\begin{proposition}[Counterexample: as an unrestricted uniqueness assertion]\label{maps:src:A.2.2}
\StatementRecord{maps:src:A.2.2}{proved}{Historical events: Claude Fable 5.2 (Anthropic), effort high, 16 November 2026: carried as background without fault found. GPT-6 Astra (OpenAI), ChatGPT Pro mode, effort Pro, S6\_09\_PAPERS, 17 November 2026: targeted proof check. Also checked in this writing. Pre-continuation writing events: GPT-6 Astra (OpenAI), ChatGPT Pro mode, effort Pro, S6\_10\_PAPERS, 18 November 2026: targeted proof check. Continuation metadata events: GPT-6 Astra (OpenAI), ChatGPT Pro mode, effort Pro, S6\_10\_PAPERS continuation, 19 November 2026: used and rechecked in this run; status-vocabulary reconciliation; metadata reconciliation only; no new mathematical check. \emph{Pre-continuation recorded proof scope.} GPT-6 Astra (OpenAI, ChatGPT Pro mode), effort Pro, 2026-11-17; Claude Fable 5.2 (Anthropic), effort high, 16 November 2026; also checked in this writing by GPT-6 Astra (OpenAI), ChatGPT Pro mode, effort Pro, 18 November 2026}{the hypotheses, notation and definitions stated here; the explicit constructions and calculations in the proof}


Agreement with the literal quotient outside a neighbourhood
of $W$, without a condition on the image of that neighbourhood,
does not determine a homotopy class $S^6\to S^3$.  All twelve
classes occur among maps satisfying that bare agreement
condition.

\end{proposition}


\begin{proof}
\begin{enumerate}[label=\arabic*.,leftmargin=*]
\item Take a closed embedded six-ball in the interior of the cusp
neighbourhood.  Homotope a given extension on a smaller ball
to be constant there, keeping its boundary and the complement
fixed; a target coordinate ball around the image of its centre
provides the contraction.  Pinch the smaller ball to
$S^6\vee D^6$ relative to its boundary and insert any based
map $\alpha:S^6\to S^3$ on the new sphere.  The resulting
global map agrees with the original map off that ball, and
represents $[e]+\alpha$ under $[S^6,S^3]=\pi_6(S^3)$.
Distinct $\alpha$ give distinct classes.  Smooth relative
approximation preserves this conclusion.  The insertion
cannot be performed within the cusp solid torus, since that
solid torus has $\pi_6=0$; this is exactly the distinction
in \cref{mapsI:src:A.2.1}.
\end{enumerate}
\end{proof}

The corrected content of \cite[Remark 8.3]{mapsEngel} is
therefore a quotient over the regular part and the two finite
fillings, followed by a replacement within its circle homotopy
class over a cusp collar.  The literal character does not
extend.  The replacement gives a canonical homotopy class
when the solid-torus image condition is retained.


\phantomsection\label{qa:III:layout:12}
\subsection{The framed three-torus and the class of \texorpdfstring{$e$}{e}}

We use $\pi_3^{\mathrm S}=\mathbb Z/24\langle\nu\rangle$, with
$e_{\mathbb R}(\nu)=1/24$, and
$\pi_6(S^3)=\mathbb Z/12\langle\nu'\rangle$, where the
stabilization of $\nu'$ is $2\nu$.  We write
$\eta_n:S^{n+1}\to S^n$ for the $(n-2)$-fold suspension of
the complex Hopf map.  The order-two classes below are
unaffected by reversal of either source or target orientation.
An orientation-preserving diffeomorphism to the standard sphere
transports its bounding stable framing to $TX\oplus\mathbb R$.
Its homotopy class does not depend on that diffeomorphism:
two oriented stable trivializations differ by a map
$S^6\to\mathrm{SO}$, and $\pi_6(\mathrm{SO})=0$
by \cite[Corollary to Theorem II and the table on p.~315]{mapsBott59}.
We use outward stabilization and the convention that the normal
frame precedes the tangential frame; a constant permutation is
corrected once to preserve orientation.

\begin{lemma}\label{maps:src:B.1.1}
\StatementRecord{maps:src:B.1.1}{proved}{Historical events: Claude Fable 5.2 (Anthropic), effort high, 16 November 2026: carried as background without fault found. GPT-6 Astra (OpenAI), ChatGPT Pro mode, effort Pro, S6\_09\_PAPERS, 17 November 2026: targeted proof check. Also checked in this writing. Pre-continuation writing events: GPT-6 Astra (OpenAI), ChatGPT Pro mode, effort Pro, S6\_10\_PAPERS, 18 November 2026: targeted proof check. Continuation metadata events: GPT-6 Astra (OpenAI), ChatGPT Pro mode, effort Pro, S6\_10\_PAPERS continuation, 19 November 2026: used and rechecked in this run; status-vocabulary reconciliation; metadata reconciliation only; no new mathematical check. \emph{Pre-continuation recorded proof scope.} GPT-6 Astra (OpenAI, ChatGPT Pro mode), effort Pro, 2026-11-17; Claude Fable 5.2 (Anthropic), effort high, 16 November 2026; also checked in this writing by GPT-6 Astra (OpenAI), ChatGPT Pro mode, effort Pro, 18 November 2026}{the hypotheses, notation and definitions stated here; the explicit constructions and calculations in the proof}


Let $F$ be a regular fibre of the original fibration.  Compare
the ambient stable framing of the oriented sphere $X$ with
the framing on $TX|_F$ consisting of four translation-invariant
vertical vector fields and an oriented lifted base frame.
The restriction of their stable difference to any coset of
$T_L\subset F$ is null-homotopic.

\end{lemma}


\begin{proof}
\begin{enumerate}[label=\arabic*.,leftmargin=*]
\item After a constant orientation adjustment the comparison is
a map $D:F\to\mathrm{SO}$, where stabilization is understood.
Its homotopy class is invariant under both monodromies.
Indeed, transport $F$ around a based loop in $B^\circ$ using
the flat real marking.  The translated vertical frame varies
continuously on $F\times[0,1]$; at the endpoint it changes by
the constant real matrix of the affine monodromy.  An oriented
base frame can be chosen periodic along the loop.  The space
of lifts of that frame to $TX$ is affine, so changing the lift
does not affect the comparison class.  The ambient stable
framing is already defined along this entire cylinder.  Thus
the comparison at its endpoints is $D$ and $D\circ A_j$
times a constant positive matrix.  Every such constant
matrix is path connected to the identity.  This proves
$[D]=A_j^*[D]$.  An affine translational part is isotopic
to the identity and has no effect.

\item The map on fundamental groups defines
$\sigma(D)\in H^1(F;\mathbb Z/2)$.  The exhaustive calculation
in \texttt{maps\_exact.py} gives
\[
 H^1(F;\mathbb Z/2)^{\langle A_1,A_2\rangle}
      =\{0,\gamma\bmod2\}.
\]
Both classes vanish on $L$, so $D|_{T_L}$ lifts to
$\mathrm{Spin}$.  Up to dimension four the latter has
$\pi_1=\pi_2=\pi_4=0$, $\pi_3=\mathbb Z$.
Obstruction theory consequently identifies maps from a
three- or four-dimensional torus to $\mathrm{Spin}$ with
$H^3$ of that torus.  The two possible lifts differ by
multiplication by the central element $-1$, which is
homotopic to the identity, so this identification has no
unrecorded lifting ambiguity.

\item The doubled stable difference $2[D]$ has zero $\sigma$,
and therefore corresponds to a monodromy-invariant class
in $H^3(F;\mathbb Z)$.  Integral Poincare duality, natural
for these determinant-one monodromies, identifies the
invariants in $H^3(F;\mathbb Z)$ with the fixed vectors
$\mathbb Zd\subset H_1(F;\mathbb Z)$.  Such a class
evaluates to zero on $T_L$, since $d\in L$:
its evaluation is the determinant with columns $d,u,w,d$.
Hence $2[D|_{T_L}]=0$ in $[T_L,\mathrm{Spin}]=\mathbb Z$.
This group has no torsion; thus $[D|_{T_L}]=0$.
The same argument applies to a coset since translation is
isotopic to the identity.
\end{enumerate}
\end{proof}

\begin{theorem}\label{maps:src:B.1.2}
\StatementRecord{maps:src:B.1.2}{proved}{Historical events: Claude Fable 5.2 (Anthropic), effort high, 16 November 2026: carried as background without fault found. Claude Fable 5.2 (Anthropic), effort high, 17 November 2026: carried as background without fault found. GPT-6 Astra (OpenAI), ChatGPT Pro mode, effort Pro, S6\_09\_PAPERS, 17 November 2026: targeted proof check. Also checked in this writing. Pre-continuation writing events: GPT-6 Astra (OpenAI), ChatGPT Pro mode, effort Pro, S6\_10\_PAPERS, 18 November 2026: targeted proof check. Continuation metadata events: GPT-6 Astra (OpenAI), ChatGPT Pro mode, effort Pro, S6\_10\_PAPERS continuation, 19 November 2026: used and rechecked in this run; status-vocabulary reconciliation; metadata reconciliation only; no new mathematical check. \emph{Pre-continuation recorded proof scope.} GPT-6 Astra (OpenAI, ChatGPT Pro mode), effort Pro, 2026-11-17; Claude Fable 5.2, effort high, 2026-11-17; Claude Fable 5.2 (Anthropic), effort high, 16 November 2026; also checked in this writing by GPT-6 Astra (OpenAI), ChatGPT Pro mode, effort Pro, 18 November 2026}{the hypotheses, notation and definitions stated here; \cref{maps:src:B.1.1}; \cite[Theorem (4.14), p.~423]{mapsAPS2}; \cite[Theorem 4.1 and Corollaries 4.2--4.3]{mapsToda52}}


For the repaired quotient of \cref{mapsI:src:A.2.1} with $|p|=1$,
\[
 [e]=6\nu'=\eta_3\eta_4\eta_5\ne0
       \quad\text{in }\pi_6(S^3)=\mathbb Z/12.
 \tag{B.1.1}\label{maps:eq:B.1.1}
\]
Its stable class is $12\nu$ and its real Adams invariant
is $1/2\in\mathbb Q/\mathbb Z$.

\end{theorem}


\begin{proof}
\begin{enumerate}[label=\arabic*.,leftmargin=*]
\item Choose a point of $Q$ over a regular base value outside
the modified cusp disc and outside the two finite root
discs.  Its entire inverse image is one coset of $T_L$.
Order a target frame as the positive real base frame followed
by the positive $\gamma$-direction.  Lifting this frame
under $de$ gives a normal three-framing.  Complete it by
the constant Lie frame of $L$ with the unique constant
orientation adjustment making the full frame positive.
This is precisely the restriction of the product frame
in Lemma \ref{maps:src:B.1.1}, up to a constant positive change of basis.
That lemma proves that the stable tangential framing of
the inverse image is the Lie framing of $T^3$.
Pontryagin--Thom therefore gives
\[
 \operatorname{st}[e]=[T^3,\mathrm{Lie}]
 = [S^1,\mathrm{Lie}]^3=\eta^3=12\nu.
 \tag{B.1.2}\label{maps:eq:B.1.2}
\]
The product equality follows by taking the product of the
three framed circles and their tubular collapse maps.
The Lie-framed circle is the nonbounding framed circle;
its class is the stabilized Hopf element $\eta$.

\item Here is an independent evaluation of the framed bordism
class, which does not use the ring relation $\eta^3=12\nu$.
Put the flat product metric on $\mathbb R^3/\mathbb Z^3$,
with the spin structure determined by its Lie framing.
Its spinors are periodic $\mathbb C^2$-valued functions.
On the Fourier mode $e^{2\pi i k\cdot x}$, $k\in\mathbb Z^3$,
the Dirac operator is the Hermitian matrix
$2\pi\sum k_j\sigma_j$, with eigenvalues
$\pm2\pi|k|$ for $k\ne0$ and a two-dimensional kernel
for $k=0$.  Thus its eta function is identically zero in
its convergent half-plane and hence at zero, while $h=2$.
The Levi-Civita connection and the connection induced
by the framing coincide, so their transgression term
vanishes.  The real formula of
\cite[Theorem (4.14), p.~423]{mapsAPS2}, with its factor
$\epsilon(1)=1/2$, gives
\[
 e_{\mathbb R}(T^3,\mathrm{Lie})
       =\frac12\frac{0+2}{2}=\frac12\pmod{\mathbb Z}.
 \tag{B.1.3}\label{maps:eq:B.1.3}
\]
This detects $12\nu$ in $\mathbb Z/24$.  The complex
$e$-invariant would be zero here; omitting the real
factor would therefore lose the desired detection.

\item Finally stabilization is injective on $\pi_6(S^3)$.
For the exact unstable input we use
\cite[Theorem 4.1 and Corollaries 4.2--4.3]{mapsToda52}:
the suspended generator is the torsion term in
$[\iota_4,\iota_4]=2\nu_4-E\nu'$, and stabilization
sends it to $2\nu$.  In the group of order twelve this
is injective.  Therefore the class with stable image
$12\nu$ is $6\nu'$.  The composite
$\eta_3\eta_4\eta_5$ has the same stable image, so the
same injection identifies it with $6\nu'$.
\end{enumerate}
\end{proof}

\subsection{The holomorphic map to the base}

\begin{theorem}\label{maps:src:B.2.1}
\StatementRecord{maps:src:B.2.1}{proved}{Historical events: Claude Fable 5.2 (Anthropic), effort high, 16 November 2026: carried as background without fault found. Claude Fable 5.2 (Anthropic), effort high, 17 November 2026: carried as background without fault found. GPT-6 Astra (OpenAI), ChatGPT Pro mode, effort Pro, S6\_09\_PAPERS, 17 November 2026: targeted proof check. Also checked in this writing. Pre-continuation writing events: GPT-6 Astra (OpenAI), ChatGPT Pro mode, effort Pro, S6\_10\_PAPERS, 18 November 2026: targeted proof check. Continuation metadata events: GPT-6 Astra (OpenAI), ChatGPT Pro mode, effort Pro, S6\_10\_PAPERS continuation, 19 November 2026: used and rechecked in this run; status-vocabulary reconciliation; metadata reconciliation only; no new mathematical check. \emph{Pre-continuation recorded proof scope.} GPT-6 Astra (OpenAI, ChatGPT Pro mode), effort Pro, 2026-11-17; Claude Fable 5.2, effort high, 2026-11-17; Claude Fable 5.2 (Anthropic), effort high, 16 November 2026; also checked in this writing by GPT-6 Astra (OpenAI), ChatGPT Pro mode, effort Pro, 18 November 2026}{the hypotheses, notation and definitions stated here; the explicit constructions and calculations in the proof}


The map $f:X\to\mathbb P^1$ of the original sphere
construction is null-homotopic.  This conclusion does not
require the choice of $e$.

\end{theorem}


\begin{proof}
\begin{enumerate}[label=\arabic*.,leftmargin=*]
\item The quotient projection $h:Q\to B$ is, after the
orientation-preserving identifications fixed in \cref{mapsI:src:A.1.3},
the weighted Hopf map
$[z_1^3:z_2^4]$, followed by a Mobius transformation
moving its two exceptional values to $0,1$.
It is the ordinary Hopf map after the self-map
\[
 P_{3,4}(z_1,z_2)=
 \frac{(z_1^3,z_2^4)}{\sqrt{|z_1|^6+|z_2|^8}}
 \quad\text{of }S^3.
 \tag{B.2.1}\label{maps:eq:B.2.1}
\]
At a value with both coordinates nonzero, its phase
map has $3\cdot4=12$ preimages, all with positive local
degree; its radial coordinate is strictly increasing.
Consequently its degree is twelve.  A degree-$k$ map
of $S^3$ is homotopic to the $k$th power map of the
topological group $S^3$.  On $\pi_6(S^3)$ the latter
is multiplication by $k$, since pointwise multiplication
and the homotopy-group addition agree.  Thus
$h_* = (\eta_2)_*\circ12=0$ on $\pi_6(S^3)$.
The construction of $e$ is over $B$, so $f=h\circ e$
as actual maps.  This proves $[f]=0$.

\item For a second proof, let $S_1,S_2$ be the reduced multiple
fibres.  The line bundles $\mathcal O_X(S_j)$ are
topologically trivial because $H^2(X;\mathbb Z)=0$.
The coordinate sections of $f^*\mathcal O(1)$ have
divisors $3S_1$ and $4S_2$ after a Mobius change of
base coordinate.  Choose a smooth frame of
$f^*\mathcal O(1)$.  Since $H^1(X;\mathbb Z)=0$, every
map $X\to\mathbb C^*$ admits a logarithm.  We can
therefore choose frames of $\mathcal O_X(S_1)$ and
$\mathcal O_X(S_2)$ whose third and fourth powers,
under the given divisor isomorphisms, equal the chosen
frame.  The canonical divisor sections then become
smooth functions $a,b$ with no common zero and
$f=[a^3:b^4]$.  Put
$G=(a,b)/\sqrt{|a|^2+|b|^2}:X\to S^3$.
Writing $r=\sqrt{|a|^2+|b|^2}>0$, the explicit homotopy
$[a^3/r^{3s}:b^4/r^{4s}]$, $0\leq s\leq1$, gives
$f\simeq\eta_2P_{3,4}G$.  It is defined everywhere because
$a$ and $b$ have no common zero.  The latter map is zero
by the same degree calculation and the exponent twelve
of $\pi_6(S^3)$.  This proof uses neither the cusp
replacement nor the computed class of $e$.
\end{enumerate}
\end{proof}

\subsection{The four workbench maps}

The exact maps, including their suspension conventions, are
needed here.  Put $\overline T=(\overline\Lambda\otimes\mathbb R)
/\overline\Lambda$ and let $\overline\pi:N_0\setminus W
\to\overline T$ be the long-period quotient.  On the punctured
cusp write $q=re^{i\phi}$.  The map $q_2$ takes $x$ to
$[r,e^{i\phi},\overline\pi(x)]$ in the unreduced suspension
of $S^1\times\overline T$, sends $W$ to its lower vertex,
and sends the complement of the cusp neighbourhood to its
upper vertex.  Collapsing the two-skeleton of
$S^1\times\overline T=T^3$ and then taking the reduced
suspension gives $H:X\to S^4$.  Collapsing the one-skeleton
of $\overline T=T^2$ and forgetting $\phi$ gives
$\widehat q_3:X\to S^3$.  Retaining only $\gamma$ and
forgetting $\phi$ gives $\widehat q_\gamma:X\to S^2$.
These formulas specify the marked maps, including the retained coordinates
and quotient identifications; arbitrary sphere quotients of the same targets
are not substituted for them.

\begin{theorem}\label{maps:src:B.3.1}
\StatementRecord{maps:src:B.3.1}{proved}{Historical events: Claude Fable 5.2 (Anthropic), effort high, 16 November 2026: carried as background without fault found. GPT-6 Astra (OpenAI), ChatGPT Pro mode, effort Pro, S6\_09\_PAPERS, 17 November 2026: targeted proof check. Also checked in this writing. Pre-continuation writing events: GPT-6 Astra (OpenAI), ChatGPT Pro mode, effort Pro, S6\_10\_PAPERS, 18 November 2026: targeted proof check. Continuation metadata events: GPT-6 Astra (OpenAI), ChatGPT Pro mode, effort Pro, S6\_10\_PAPERS continuation, 19 November 2026: used and rechecked in this run; status-vocabulary reconciliation; metadata reconciliation only; no new mathematical check. \emph{Pre-continuation recorded proof scope.} GPT-6 Astra (OpenAI, ChatGPT Pro mode), effort Pro, 2026-11-17; Claude Fable 5.2 (Anthropic), effort high, 16 November 2026; also checked in this writing by GPT-6 Astra (OpenAI), ChatGPT Pro mode, effort Pro, 18 November 2026}{the hypotheses, notation and definitions stated here; \cref{maps:eq:B.3.2}; \cref{maps:eq:B.3.3}}


For the exact retained-angle map just defined,
\[
 [H]=\eta_4\eta_5\ne0\in\pi_6(S^4)=\mathbb Z/2.
 \tag{B.3.1}\label{maps:eq:B.3.1}
\]
Its regular inverse image is the native vanishing
two-torus with its Lie stable framing and Arf invariant one.

\end{theorem}


\begin{proof}
\begin{enumerate}[label=\arabic*.,leftmargin=*]
\item Take $u_0=(-1/4,1/4)$, $q_0=r_0e^{i\pi}$, and
$p_0=B_0u_0=(1/4,1/4)$.  For $r_0$ sufficiently small
the full inverse image of the corresponding open-cell
value lies in the original chart of \cref{mapsI:src:A.1.2}.  To see this
without losing components, a lift of any point of that
inverse image has long coordinate $u_0+n$ for a unique
$n\in\mathbb Z^2$; applying the actual period deck
transformation $-n$ changes it to $u_0$.  Distinct lifts
with this coordinate cannot be deck equivalent.
Put $h_j(q)=\exp(-2\pi i(C(q)u_0)_j)$ and make the
holomorphic coordinate change
\[
 (w_0,w_1,w_2)
   =(z_0/(h_1h_2),h_1z_1,h_2z_2).
\]
It preserves the product $q$.  The full inverse image is
\[
 w(a,b)=
 (r_0^{1/2}e^{i(\pi-a-b)},
  r_0^{1/4}e^{ia},r_0^{1/4}e^{ib}),
 \quad a,b\in\mathbb R/2\pi\mathbb Z.
 \tag{B.3.2}\label{maps:eq:B.3.2}
\]
All coordinates tend uniformly to zero as $r_0\to0$.
The period quotient is injective on a neighbourhood of
the chart origin, so this torus is contained in a
contractible coordinate ball of $X$.
The map $H$ is smooth on an open neighbourhood of this
torus, although its suspension presentation is only
continuous elsewhere.  Choose a smaller open-cell target
ball whose full inverse image lies in that neighbourhood;
compactness permits this because the displayed torus is
the full inverse image of its centre.  Relative smooth
approximation, uniformly small off that neighbourhood,
then gives a smooth representative with exactly the same
inverse image and differential there.  Pontryagin--Thom
therefore applies to the framing computed below.

\item We compute its framing before discarding the bounded
period terms.  Write
\[
 \ell=\log|q|,\quad
 R(q)=-2\pi\operatorname{Im}C(q),\quad
 A(q)=\ell B_0+R(q),\quad
 v=(\log|w_1|,\log|w_2|)^t.
\]
The exact coordinate identity is
$v=\ell B_0u+R(q)(u-u_0)$, and on \eqref{maps:eq:B.3.2} it gives
\[
 du=A(q_0)^{-1}(dv-p_0\,d\ell).
 \tag{B.3.3}\label{maps:eq:B.3.3}
\]
Thus no derivative of a bounded correction has been
ignored.  Let $L_j$ and $Q_j$ denote the real tangent
vectors $w_j\partial_{w_j}$ and $iw_j\partial_{w_j}$
in the real identification $T\mathbb C^3=\mathbb C^3$.
With $A_0=A(q_0)$, dual lifts of the target coordinate
order $(\ell,\phi,u_1,u_2)$ are
\[
 N_\ell=\tfrac12L_0+\tfrac14L_1+\tfrac14L_2,
 \qquad N_\phi=Q_0,
\]
\[
 N_i=-((A_0)_{1i}+(A_0)_{2i})L_0
            +(A_0)_{1i}L_1+(A_0)_{2i}L_2,
 \quad i=1,2.
 \tag{B.3.4}\label{maps:eq:B.3.4}
\]
The tangent vectors are $T_a=-Q_0+Q_1$ and
$T_b=-Q_0+Q_2$.  Formula \eqref{maps:eq:B.3.3} checks that the
four normal vectors map to the four standard coordinate
vectors and both tangent vectors map to zero.
The full determinant is
$-\det(A_0)|q_0|^2<0$, so the positive tangent order is
$(-T_a,T_b)$.  The inequality holds for small $r_0$
because $A_0$ is joined to $\ell_0B_0$ through
$\ell_0B_0+sR(q_0)$, all of which are invertible.

\item Every displayed vector is a constant linear combination
of $L_j,Q_j$.  Its dependence on $(a,b)$ is therefore
exactly
\[
 D(a,b)=\operatorname{diag}(e^{-i(a+b)},e^{ia},e^{ib})
            \in\mathrm{SU}(3).
\]
This map is null-homotopic.  An explicit periodic
contraction uses
\[
 Q_s(\theta)=
 \begin{pmatrix}
 \sin(\pi s/2)e^{-i\theta}&-\cos(\pi s/2)\\
 \cos(\pi s/2)&\sin(\pi s/2)e^{i\theta}
 \end{pmatrix},
 \qquad K_s(\theta)=Q_s(\theta)Q_s(0)^*.
\]
Each matrix lies in $\mathrm{SU}(2)$,
$K_0=I$, and $K_1(\theta)=\operatorname{diag}
(e^{-i\theta},e^{i\theta})$.  The product of its
embeddings in coordinate planes $(0,1)$ and $(0,2)$,
evaluated at $a$ and $b$, contracts $D$.
Consequently the induced stable tangential frame of
\eqref{maps:eq:B.3.2} is the constant Lie frame.  The ambient framing
comparison on the containing ball is null-homotopic,
so the same conclusion holds in the framed sphere.

\item For the induced spin structure the two coordinate
circles are nonbounding.  Its quadratic form on
$H_1(T^2;\mathbb F_2)$ is
$q(xa+yb)=x+xy+y$, and its Arf invariant is
$q(a)q(b)=1$.  Hence its framed class is
$\eta^2\in\pi_2^{\mathrm S}$.  Freudenthal suspension
is an isomorphism for $\pi_6(S^4)$, so the unstable
class is exactly \eqref{maps:eq:B.3.1}.
\end{enumerate}
\end{proof}

\begin{theorem}\label{maps:src:B.3.2}
\StatementRecord{maps:src:B.3.2}{proved}{Historical events: Claude Fable 5.2 (Anthropic), effort high, 16 November 2026: carried as background without fault found. GPT-6 Astra (OpenAI), ChatGPT Pro mode, effort Pro, S6\_09\_PAPERS, 17 November 2026: targeted proof check. Also checked in this writing. Pre-continuation writing events: GPT-6 Astra (OpenAI), ChatGPT Pro mode, effort Pro, S6\_10\_PAPERS, 18 November 2026: targeted proof check. Continuation metadata events: GPT-6 Astra (OpenAI), ChatGPT Pro mode, effort Pro, S6\_10\_PAPERS continuation, 19 November 2026: used and rechecked in this run; status-vocabulary reconciliation; metadata reconciliation only; no new mathematical check. \emph{Pre-continuation recorded proof scope.} GPT-6 Astra (OpenAI, ChatGPT Pro mode), effort Pro, 2026-11-17; Claude Fable 5.2 (Anthropic), effort high, 16 November 2026; also checked in this writing by GPT-6 Astra (OpenAI), ChatGPT Pro mode, effort Pro, 18 November 2026}{the hypotheses, notation and definitions stated here; the explicit constructions and calculations in the proof}


The angle-forgetting workbench maps satisfy
$[\widehat q_3]=0$ and $[\widehat q_\gamma]=0$.

\end{theorem}


\begin{proof}
\begin{enumerate}[label=\arabic*.,leftmargin=*]
\item Both components of $\overline\pi$ extend as integral
$H^1$ classes over $N_0$, by (A.1.5).  Choose a map
$\pi':N_0\to T^2$ representing them and equal to
$\overline\pi$ on an outer collar; the construction
and relative uniqueness are as in \cref{mapsI:src:A.2.1}.  On the
punctured cusp the quotient $\overline\pi-\pi'$ has
a lift to $\mathbb R^2$ that is zero on the collar.
Indeed the punctured cusp retracts onto its boundary
mapping torus, and the difference is zero in relative
$H^1$.  Multiplying that lift by a homotopy parameter
gives a homotopy from $\overline\pi$ to $\pi'$.
After passage to either suspension target this
homotopy is continuous at $W$, uniformly in the
homotopy parameter: the whole compact torus of
angles is collapsed as $r\to0$.  It is also constant
on the outer boundary and extends by the basepoint
outside $N_0$.

\item After this homotopy the three-sphere map factors as
\[
 X\longrightarrow
 \frac{[0,\epsilon]\times T^2}{\{\epsilon\}\times T^2}
 \longrightarrow\Sigma(T^2/T^{2,(1)}),
 \qquad x\longmapsto[r(x),\pi'(x)]
\]
on $N_0$, with the complement mapping to the vertex.
The middle space is a cone and is contractible.
For the two-sphere map use the first component
$\theta':N_0\to S^1$ and the cone on $S^1$.
These give actual null-homotopies of the two maps.
This argument cannot discard the retained base
angle in $H$, whose intermediate space is instead
$(D^2\times T^2)/(\partial D^2\times T^2)$.
\end{enumerate}
\end{proof}

\begin{theorem}\label{maps:src:B.3.3}
\StatementRecord{maps:src:B.3.3}{proved}{Historical events: Claude Fable 5.2 (Anthropic), effort high, 16 November 2026: carried as background without fault found. GPT-6 Astra (OpenAI), ChatGPT Pro mode, effort Pro, S6\_09\_PAPERS, 17 November 2026: targeted proof check. Also checked in this writing. Pre-continuation writing events: GPT-6 Astra (OpenAI), ChatGPT Pro mode, effort Pro, S6\_10\_PAPERS, 18 November 2026: targeted proof check. Continuation metadata events: GPT-6 Astra (OpenAI), ChatGPT Pro mode, effort Pro, S6\_10\_PAPERS continuation, 19 November 2026: used and rechecked in this run; status-vocabulary reconciliation; metadata reconciliation only; no new mathematical check. \emph{Pre-continuation recorded proof scope.} GPT-6 Astra (OpenAI, ChatGPT Pro mode), effort Pro, 2026-11-17; Claude Fable 5.2 (Anthropic), effort high, 16 November 2026; also checked in this writing by GPT-6 Astra (OpenAI), ChatGPT Pro mode, effort Pro, 18 November 2026}{the hypotheses, notation and definitions stated here; \cref{maps:src:B.3.2}; \cref{maps:src:B.1.2}}


Let $C=M\cup\operatorname{Cone}(\partial M)$ and let
$q_1:X\to C$ be the workbench map retaining the
Seifert quotient away from the cusp and collapsing
$W$ to the cone point.  Under the cofiber splitting
$C\simeq S^3\vee S^2$ determined by the cusp core,
\[
 [q_1]=(j_3)_*(6\nu'),
 \tag{B.3.5}\label{maps:eq:B.3.5}
\]
with no $S^2$ or mixed Whitehead component.

\end{theorem}


\begin{proof}
\begin{enumerate}[label=\arabic*.,leftmargin=*]
\item Write $Q=M\cup(D^2\times S^1)$ and let $K$ be the
core $\{0\}\times S^1$ of its cusp solid torus.
Collapsing this core turns the solid torus into
the cone on its boundary, by its radial coordinate.
Thus $C=Q/K$ as a quotient space.  The same
homotopy used in \ref{maps:src:B.3.2}, now retaining the base angle,
shows $q_1\simeq\kappa e$, where $\kappa:Q\to Q/K$
is the quotient map.  The inclusion $K=S^1\to S^3$
is null-homotopic, and its cofiber is
$S^3\vee\Sigma S^1$.  We choose this equivalence
to restrict on $Q$ to $j_3$.  A null-homotopy of
the circle exists since $\pi_1(S^3)=0$; its
ambiguity is $\pi_2(S^3)=0$.  Hence the indicated
splitting is fixed in the only respect used here,
and \ref{maps:src:B.1.2} proves \eqref{maps:eq:B.3.5}.
\end{enumerate}
\end{proof}

\subsection{The other named sphere fibrations}

\begin{theorem}\label{maps:src:B.4.1}
\StatementRecord{maps:src:B.4.1}{proved}{Historical events: Claude Fable 5.2 (Anthropic), effort high, 16 November 2026: carried as background without fault found. GPT-6 Astra (OpenAI), ChatGPT Pro mode, effort Pro, S6\_09\_PAPERS, 17 November 2026: targeted proof check. Also checked in this writing. Pre-continuation writing events: GPT-6 Astra (OpenAI), ChatGPT Pro mode, effort Pro, S6\_10\_PAPERS, 18 November 2026: targeted proof check. Continuation metadata events: GPT-6 Astra (OpenAI), ChatGPT Pro mode, effort Pro, S6\_10\_PAPERS continuation, 19 November 2026: used and rechecked in this run; status-vocabulary reconciliation; metadata reconciliation only; no new mathematical check. \emph{Pre-continuation recorded proof scope.} GPT-6 Astra (OpenAI, ChatGPT Pro mode), effort Pro, 2026-11-17; Claude Fable 5.2 (Anthropic), effort high, 16 November 2026; also checked in this writing by GPT-6 Astra (OpenAI), ChatGPT Pro mode, effort Pro, 18 November 2026}{the hypotheses, notation and definitions stated here; \cite[Lemma 2.1 and equations (2.11)--(2.12)]{mapsViaLog}; \cref{maps:src:B.1.1}; \cref{maps:src:B.1.2}; \cite[Proposition 4.2]{mapsVia2}; \cite[Proposition 5.4]{mapsVia2}; \cite[Lemma 3.10 and Proposition 3.11]{mapsLR}; \cite[Proposition 4.8]{mapsLR}; \cite[Propositions 3.12--3.13 and Corollary 3.16]{mapsLR}; \cref{maps:eq:B.4.1}; \cref{maps:eq:B.4.3}; \cref{maps:src:B.3.1}; \cref{maps:src:B.3.2}}


For the sphere members of Viaclovsky's two-parameter family,
the Liu--Rollenske family, and every $X_{3k,4l}$ with
$\gcd(3k,4l)=1$, the primitive invariant-circle quotient
admits a canonical replacement over the degeneration
neighbourhoods as in \cref{mapsI:src:A.2.1}.  Its target is $S^3$ and
its class is $6\nu'$.  In each case the map to
$\mathbb P^1$ is null-homotopic.  The retained-angle
map at the primitive rank-two cusp is $\eta_4\eta_5$,
and the two corresponding angle-forgetting maps are
null-homotopic.  Thus the quotient class $[e]$ does
not distinguish these fibrations.

\end{theorem}


\begin{proof}
\begin{enumerate}[label=\arabic*.,leftmargin=*]
\item We supply the integral lattices, the local extension
comparisons, and the target calculation; a list of
Betti numbers would not imply this statement.

\item For $X_{3k,4l}$ the regular homological representation
is still $(A_1,A_2,M_0)$ of \cref{mapsI:src:A.1.1}.  We use \cref{push:src:E.1.1} and \cref{push:src:E.1.2} with
its original finite-root convention.  In Viaclovsky's
marking $(\delta,a,b,d)$ the comparison is
\[
 P_V=\begin{pmatrix}-2&1&0&1\\-6&0&-1&0\\0&-1&0&0\\-1&0&0&0\end{pmatrix}.
\]
Substitution in the matrices of
\cite[Lemma 2.1 and equations (2.11)--(2.12)]{mapsViaLog}
gives $V_iP_V=P_VA_{i+1}$, $\det P_V=-1$,
$P_Vd=\delta$, $P_Vv_1=-u_0$, and
$P_V(-v_2)=\delta-u_1$.  This checks the step of
\cref{push:src:E.1.1} and \cref{push:src:E.1.2} that we use, without redoing its analytic
classification.  With $m=3k,n=4l$, the actual
finite translations are the vectors in \cref{push:src:E.1.1} and \cref{push:src:E.1.2} and
the circle transgression is
$b=nN_0+mN_1=\pm1$ for a sphere member.
The two target fillings have multiplicities $m,n$;
their genus-one meridian determinant is $b$.
Thus their target is $S^3$.  Its map to the base
is a weighted Hopf map of degree product $mn=12kl$.
The common lattices and their mod-two counterparts
are those in \ref{maps:src:B.1.1}, and all the arguments of
\cref{mapsI:src:A.2.1} and \ref{maps:src:B.1.2} apply without change.

\item For the two-parameter family use the homology basis
$(\widehat\psi,\widehat u,\widehat w,\widehat\delta)$
of \cref{push:src:E.5.1} and \cref{push:src:E.5.2} and \cref{push:src:E.5.3}.  Its actual four matrices are
\[
 H_* =\begin{pmatrix}1&0&0&0\\-2&0&-1&0\\0&1&0&0\\1&1&0&1\end{pmatrix},
 \qquad
 H_1 =\begin{pmatrix}1&0&0&0\\0&1&0&0\\-2&-1&1&0\\-2&-1&0&1\end{pmatrix},
\]
\[
 H_2 =\begin{pmatrix}1&0&0&0\\0&1&1&0\\0&0&1&0\\0&0&0&1\end{pmatrix},
 \qquad
 H_3 =\begin{pmatrix}1&0&0&0\\0&1&0&0\\0&-1&1&0\\1&0&0&1\end{pmatrix}.
\]
They are inverse transposes of the matrices in
\cite[Proposition 4.2]{mapsVia2}.  The direct
calculation gives
\[
 \Lambda^\Gamma=\mathbb Z\widehat\delta,
 \quad (\Lambda^*)^\Gamma=\mathbb Z\psi,
 \quad (\Lambda^*\otimes\mathbb F_2)^\Gamma
          =\mathbb F_2\psi,
 \quad L=\ker\psi.
 \tag{B.4.1}\label{maps:eq:B.4.1}
\]
The finite root action has order four and translation
$\ell/4$, where $\ell=(1,-1,-1,0)^t$; hence
$\psi(\ell)=1$.  The three primitive vanishing
lattices are
\[
 K_1=\mathbb Z(\widehat w+\widehat\delta),\qquad
 K_2=\mathbb Z\widehat u,\qquad
 K_3=\mathbb Z\langle\widehat w,\widehat\delta\rangle.
 \tag{B.4.2}\label{maps:eq:B.4.2}
\]
They all lie in $L$.  To check that the needed
class really extends, the two rank-one local
models in the source are the quotient of an
infinite chain of ruled elliptic surfaces, with
the adjacent sections identified by a translation.
Their nearby fundamental group surjects to the
central one with precisely $K_i$ killed; the
base meridian bounds the extending section.
This follows directly by filling the one
vanishing circle with its nodal thimble, and
retaining the other three circle generators.
It is also the degree-one part of the source's
local-specialization proposition, labelled
\texttt{mum}.  At the rank-two cusp the same
assertion is (A.1.5).  Thus $\psi$ extends
over every reduced filling, and its boundary
restriction is injective on $H^1$.  \cref{mapsI:src:A.2.1}
constructs and compares the replacements.

\item There is a single multiple target circle of
multiplicity four.  The other three circle
fillings are ordinary solid tori, with no
meridian shift because the chosen section
extends there.  The finite logarithmic
shift has primitive circle component $1/4$.
Equivalently its integral transgression is
$d_2(4\psi)=1$, as computed directly in
\cite[Proposition 5.4]{mapsVia2}.
Hence the genus-one target has meridian
determinant one and is $S^3$.

\item For Liu--Rollenske use their common homology
basis $(e_1,e_2,e_3,e_4)$.  The two finite
monodromies are
\[
 R_0=\begin{pmatrix}-1&-1&-1&0\\2&1&1&0\\0&0&1&0\\0&0&0&1\end{pmatrix},
 \qquad
 R_1=\begin{pmatrix}1&1&0&1\\-2&0&-1&-1\\0&-1&1&-1\\0&0&0&1\end{pmatrix}.
\]
Their source uses $R_\infty=(R_1R_0)^{-1}$.
They have respective orders four and six,
but the fibre over $1$ is reduced: finite
monodromy does not mean it is a smooth
multiple fibre.  The common invariants are
\[
 \Lambda^\Gamma=\mathbb Z(-e_1+2e_3),
 \quad(\Lambda^*)^\Gamma=\mathbb Z e_4^*,
 \quad(\Lambda^*\otimes\mathbb F_2)^\Gamma
            =\mathbb F_2 e_4^*.
 \tag{B.4.3}\label{maps:eq:B.4.3}
\]
Here $\delta=-e_1+2e_3$ and $\xi=e_4^*$,
so again the fixed vector belongs to
$L=\ker\xi$.
The local product identification of
\cite[Lemma 3.10 and Proposition 3.11]{mapsLR},
$Y_{\Delta_1}\cong S_{\Delta_1}\times T^2$,
preserves the two circle classes $f_3,f_4$,
with $\xi=f_4^*$.  Consequently $\xi$
extends over this reduced type-II fibre
and is surjected onto by the nearby
fundamental group.  At infinity the
primitive vanishing lattice is
$\mathbb Z\langle h_1,h_3\rangle$ and
$\xi=h_4^*$, so (A.1.5) gives its
extension there.  The only logarithmic
shift is $a=e_4/4$ over $0$.
Its primitive transgression is
$d_2(4\xi)=1$ by \cite[Proposition 4.8]{mapsLR}.
Thus the same one-exceptional-circle
target is $S^3$.  These are applications
of the exact local product and integral
monodromy statements in
\cite[Propositions 3.12--3.13 and Corollary 3.16]{mapsLR},
with the comparison of classes written
explicitly; rational invariant cycles
alone would not suffice.

\item For both one-multiple-fibre families,
Lemma \ref{maps:src:B.1.1} applies verbatim with
$(\gamma,d)$ replaced by $(\psi,\widehat\delta)$
or $(\xi,\delta)$.  Equations \eqref{maps:eq:B.4.1}
and \eqref{maps:eq:B.4.3} are all of its hypotheses.
The chosen regular inverse image of
the quotient map is a single $T_L$,
because every replacement is over its
own degeneration disc.  Therefore
Theorem \ref{maps:src:B.1.2} gives $[e]=6\nu'$ in
each family.  The target projection
is the weighted Hopf map of weights
$(4,1)$, whose power-map factor has
degree four; it sends $6\nu'$ to zero.
For $X_{3k,4l}$ the degree is $12kl$,
which also sends it to zero.

\item Finally each named family has a
primitive hexagonal rank-two cusp.
The proof of \ref{maps:src:B.3.1} is local in its
normal-crossing chart, independent
of the other fibres and of the
bounded period matrix.  It gives
$H=\eta_4\eta_5$ with the same target
coordinate convention.  The two
long-period classes extend over
that cusp, so \ref{maps:src:B.3.2} proves the
two angle-forgetting assertions.
All kernel and mod-two calculations
above are executed in
\texttt{code/maps\_exact.py}.
\end{enumerate}
\end{proof}

\subsection{A cell presentation from the three fillings}

We give a topological cell construction, retaining the actual
boundary maps before making cellular approximations.  The
smaller cellular complex used for the integer calculation is
homotopy equivalent to this construction; its cell counts are
not asserted to be the counts of a smooth handle decomposition.
Throughout this subsection $p=-1$, so the finite twists are
$v_1=g+2u-4w$ and $v_2=-g-3u+3w$ and the cusp gluing has
zero translation drift.

For a matrix $A$ on $F=\mathbb R^4/\Lambda$, our mapping-torus
convention is
\[
 M_A=(F\times[0,1])/((x,1)\sim(A^{-1}x,0)).
 \tag{B.5.1}\label{maps:eq:B.5.1}
\]
The positive interval loop $\rho$ satisfies
$\rho t_\lambda\rho^{-1}=t_{A\lambda}$.  This fixes the
inverse in \eqref{maps:eq:B.5.1} and agrees with the clockwise meridians
used in A.1.

\begin{theorem}\label{maps:src:B.5.1}
\StatementRecord{maps:src:B.5.1}{proved}{Historical events: Claude Fable 5.2 (Anthropic), effort high, 16 November 2026: carried as background without fault found. Claude Opus 5.5 (Anthropic), effort max, sweep 1, scope s6, 18 November 2026: computation or formal certificate. GPT-6 Astra (OpenAI), ChatGPT Pro mode, effort Pro, S6\_09\_PAPERS, 17 November 2026: targeted proof check. Also checked in this writing. Pre-continuation writing events: GPT-6 Astra (OpenAI), ChatGPT Pro mode, effort Pro, S6\_10\_PAPERS, 18 November 2026: targeted proof check. Continuation metadata events: GPT-6 Astra (OpenAI), ChatGPT Pro mode, effort Pro, S6\_10\_PAPERS continuation, 19 November 2026: used and rechecked in this run; status-vocabulary reconciliation; metadata reconciliation only; no new mathematical check. \emph{Pre-continuation recorded proof scope.} GPT-6 Astra (OpenAI, ChatGPT Pro mode), effort Pro, 2026-11-17; Claude Fable 5.2 (Anthropic), effort high, 16 November 2026; also checked in this writing by GPT-6 Astra (OpenAI), ChatGPT Pro mode, effort Pro, 18 November 2026}{the hypotheses, notation and definitions stated here; \cref{maps:eq:B.5.3}; \cite{mapsAC26}; \cref{maps:eq:B.5.1}; \cref{maps:eq:B.5.6}; \cite[Theorem 6.1, p.~77]{mapsWaldhausen68}; \cref{maps:eq:B.5.4}; \cref{maps:eq:B.5.8}; \cite[Propositions 4G.1--4G.2]{mapsHatcher02}}


The underlying topological manifold of $X$ is obtained from
the flat $T^4$-bundle over a ribbon bouquet of two circles,
with monodromies $A_1,A_2$, and the three spaces $S_1,S_2,W$
by the three double mapping cylinders specified below.
Their maps admit finite polyhedral cell structures.  They
therefore give a CW structure adapted to the pair of pants
and the three special fibres, with all attaching maps
specified by integral affine formulae and primitive circle
quotients.  Replacing the maps by cellular maps on the smaller
torus cell structures gives a finite CW homotopy model with
cell counts
\[
 (c_0,c_1,c_2,c_3,c_4,c_5,c_6)
       =(5,20,45,56,42,17,3).
 \tag{B.5.2}\label{maps:eq:B.5.2}
\]

\end{theorem}


\begin{proof}
\begin{enumerate}[label=\arabic*.,leftmargin=*]
\item Give a bouquet with oriented loops $\rho_1,\rho_2$ the ribbon
order in which the two ends of each loop are consecutive.
Its oriented thickening is a pair of pants.  Collapse the
widths of its two bands to the edges and the remaining
vertex disc to the vertex.  The induced boundary maps are
piecewise linear paths with words
$\rho_1,\rho_2,(\rho_1\rho_2)^{-1}$; intervals between
successive letters are mapped to the vertex.  The thickening
is the mapping cylinder of this boundary map.  This assertion
can be checked on each band, where the cylinder consists of
the two half-bands, and on the vertex disc, where the cones
on the intervening boundary intervals are its sectors.
Lift this construction to the flat $F$-bundle.  Write $U_*$
for the resulting bundle over the bouquet.  Its boundary
spaces are $M_{A_1},M_{A_2},M_{M_0}$.  Their maps
$b_j:M_{A_j}\to U_*$ and $b_0:M_{M_0}\to U_*$ are the
lifted ribbon boundary maps.  They are the identity on a
marked fibre and send their interval loops respectively to
$\rho_1,\rho_2,\rho_2^{-1}\rho_1^{-1}$.
The equality $M_0=(A_1A_2)^{-1}$ verifies the last map
including its fibre monodromy.

\item At a finite fibre put
\[
 h_j(x)=A_jx+v_j/m_j,\qquad S_j=F/\langle h_j\rangle,
 \qquad
 r_j([x,\theta])=[x+\theta v_j/m_j].
 \tag{B.5.3}\label{maps:eq:B.5.3}
\]
The last formula is a well-defined map $M_{A_j}\to S_j$:
at $\theta=1$ its value agrees with $[A_j^{-1}x]$, because
$h_j(A_j^{-1}x)=x+v_j/m_j$.  The finite neighbourhood is
the associated disc bundle
$(F\times D^2)/\langle(x,s)\mapsto(h_jx,e^{-2\pi i/m_j}s)\rangle$.
Its boundary point corresponding to $[x,\theta]$ is
\[
 [x+\theta v_j/m_j,\ e^{-2\pi i\theta/m_j}].
\]
Radial collapse is exactly \eqref{maps:eq:B.5.3}; its mapping cylinder is
the disc bundle.  In particular this is a statement about
the underlying smooth neighbourhood and does not replace
the scheme-theoretic fibre $m_jS_j$ by a reduced fibre in
the holomorphic map.  The logarithmic gluing translation is
$\log(s)\Pi v_j/(2\pi i)$.  Its clockwise drift is precisely
$v_j/m_j$.  Subtracting that displayed drift leaves a
periodic translation with a lift to $\mathbb R^4$; linear
interpolation of this lift straightens it on a collar.
Thus the actual gluing is identified with \eqref{maps:eq:B.5.3}, without
discarding a loop of translations of nonzero lattice drift.
This also corrects the first paragraph on page 42 of
\cite{mapsAC26}: multiplicity does not prevent this
neighbourhood from being a smooth normal disc bundle.
We make no assertion of a holomorphic tubular projection.
The meridian relation proved in that paragraph remains
valid; its meridian is the normal circle winding $m_j$
times around the base, rather than a once-around base
loop.

\item Here is the complete cusp boundary map.  Put
$a=(a_1,a_2)\in\mathbb R^2/\mathbb Z^2$ for the $(g,u)$
coordinates and $b\in T_f=\mathbb R^2/\mathbb Z^2$ for
the $(w,d)$ coordinates.  Let $\mathcal H$ be the hexagonal
honeycomb dual to the $A_2$ triangulation.  For $B_0a$ in
the interior of the hexagon centred at $v\in\mathbb Z^2$,
write $v(a)=v$.  The central fibre has the explicit quotient
description
\[
 W_{\mathcal H}=(T_B\times T_f)/\sim,
 \qquad T_B=\mathbb R^2/\mathbb Z^2,
 \tag{B.5.4}\label{maps:eq:B.5.4}
\]
where over a honeycomb edge separating centres $v,v'$ one
collapses the primitive circle on $v'-v$, and over each
honeycomb vertex one collapses all of $T_f$.  No circle is
collapsed over a hexagon interior.  The three oriented edge
directions can be chosen as
\[
 n_1=w,\qquad n_2=d-w,\qquad n_3=-d.
 \tag{B.5.5}\label{maps:eq:B.5.5}
\]
The quotient is homeomorphic to the actual $W$: the positive
part of each toric component is its hexagon, the phase
stabilizers are exactly these primitive circles, and the
period action identifies hexagons by $B_0\mathbb Z^2$.
The constant phase cocycle is removed by the linear phase
gauge coming from $\operatorname{Re}C(0)a$.

\item With the interval convention \eqref{maps:eq:B.5.1}, the cusp attachment is
\[
 r_0([a,b,\theta])
   =[a,b-\theta(B_0a-v(a))]\ \in W_{\mathcal H}.
 \tag{B.5.6}\label{maps:eq:B.5.6}
\]
This formula includes its seam data.  On crossing an edge
the two expressions differ by $\theta(v-v')$, which lies
in the circle collapsed on that edge.  At a vertex every
phase is collapsed.  Replacing $a$ by $a+\lambda$ changes
$v(a)$ by $B_0\lambda$, so the expression is periodic.
At $\theta=1$, the integer vector $v(a)$ disappears modulo
$\mathbb Z^2$, leaving $[a,b-B_0a]$; this is the endpoint
identification by $M_0^{-1}$.  Hence \eqref{maps:eq:B.5.6} is continuous
on the entire mapping torus, not only on its open strata.

\item We justify that the cusp neighbourhood is its actual
mapping cylinder.  Straighten $C(q)$ to $C_0=C(0)$ by
the map on logarithmic coordinates
\[
 \xi\longmapsto
 \xi-(C(q)-C_0)(\operatorname{Im}Z_q)^{-1}
                       \operatorname{Im}\xi.
 \tag{B.5.7}\label{maps:eq:B.5.7}
\]
Writing $\xi=Z_qa+b$ proves that this sends $Z_qa+b$ to
$Z_q^0a+b$, fixes the real lattice, and conjugates the
period actions.  In a bounded-position toric chart the
correction multiplying $x$ tends to one as $q\to0$;
the same holds for its inverse.  Thus it extends by the
identity on $W$ to a homeomorphism of the closed cusp tube.

\item Let $A$ be the nonnegative real part of the frozen toric
cover over $|q|\leq\eta$, and $A_0$ its central part.
The local orthant descriptions make $A$ a three-manifold
with boundary; the corners may be rounded.  Its interior
has coordinates $(y,\rho)\in\mathbb R^2\times(0,\eta)$,
$y=\log|x|/\log\rho$.  The positive period action is
$y\mapsto y+B_\rho^+\lambda$, where
$B_\rho^+=B_0-2\pi\operatorname{Im}C_0/\log\rho$.
For small $\eta$ these matrices are invertible.  Thus the
interior of the positive quotient $Q^+$ is explicitly
$T^2\times(0,\eta)$, using $a=(B_\rho^+)^{-1}y$.
The bounded fundamental parallelogram meets only finitely
many of the toric bounded-position charts, including at
$\rho=0$, so $Q^+$ is compact.  Its two boundary components
are tori.  The inner one is the honeycomb quotient, and
the outer one is the positive regular-fibre torus.

\item Both boundary inclusions induce the identity on the period
deck group $\mathbb Z^2$.  Indeed $A_0$ is the honeycomb
plane and hence contractible; $A$ is contractible because
its interior is contractible and pushing a boundary collar
inward gives a homotopy equivalence with that interior.
The quotient is irreducible.  An embedded locally flat
two-sphere lies in its interior $T^2\times(0,\eta)$.
Lift it to $\mathbb R^2\times\mathbb R$.  Its bounded
ball is disjoint from all its nonzero lattice translates:
disjoint boundary spheres would otherwise give nested
balls, impossible for translates of equal finite volume.
The ball therefore descends.  This also proves
irreducibility of the product target $T^2\times I$.
The boundary tori are incompressible because their
fundamental groups inject.

\item There is a map of pairs $Q^+\to T^2\times I$ inducing
the identity on $\pi_1$ and restricting to the prescribed
homeomorphisms on the two boundaries.  For its first
coordinate, extend the classifying map to $T^2=K(\mathbb
Z^2,1)$ and use the homotopy extension property on the
boundary; use a collar height for the second coordinate.
Apply \cite[Theorem 6.1, p.~77]{mapsWaldhausen68} to this
map.  The source is sufficiently large, since a boundary
torus pushed into its interior is incompressible.  The
product target has the same property, and all the other
hypotheses were just verified.  The homotopy is fixed on
the boundary.  Its alternative of mapping the connected
source into the boundary is impossible, since the two
source boundaries map to distinct target components.
It therefore ends at a covering map.  The induced
fundamental-group isomorphism makes that covering
one-sheeted.  We obtain a product homeomorphism relative
to both prescribed boundary markings.

\item Lift this product to $A$ and spread it by compact phases.
The full period action on a phase $(\phi,\vartheta)$ is
$(\phi,\vartheta)\mapsto
(u_\lambda\phi\vartheta^{B_0\lambda},\vartheta)$;
it does not depend on the positive radial point.  The
lifted product is therefore equivariant for the full
action.  Phase stabilizers occur only on $A_0$ and are
exactly those of \eqref{maps:eq:B.5.4}.  Consequently the resulting
quotient is the mapping cylinder of the boundary phase
map.  On the component with centre $v$, the $q$-phase
acts on $T_f$ by $-v$ times that phase, since $(v,1)$ is
the stabilizer ray.  Substituting the clockwise interval
phase and the marked period $-\theta B_0a$ gives precisely
\eqref{maps:eq:B.5.6}.

\item Gluing the three mapping cylinders of $r_j$ to the
mapping cylinder of the ribbon boundary maps gives
\[
 U_*\ \cup\ \coprod_{j=0}^2(M_j\times[0,1])
       \ \cup\ (W\sqcup S_1\sqcup S_2),
 \tag{B.5.8}\label{maps:eq:B.5.8}
\]
with $(x,0)\sim b_j(x)$ and $(x,1)\sim r_j(x)$.
The preceding constructions identify \eqref{maps:eq:B.5.8}
homeomorphically with $X$.  The comparison need not
preserve the radial value of $f$ inside the cusp, but
it preserves the special strata and all three outer
collars; this is the sense of adaptation used here.

\item For a finite cell construction use rational hexagons,
unit cubes for tori, and the affine quotient \eqref{maps:eq:B.5.3}.
For \eqref{maps:eq:B.5.6} use toric phase coordinates
$b'=b-\theta B_0a$.  In these coordinates its expression
on a hexagon is $[a,b'+\theta v]$, and its face
identifications are affine:
$a\mapsto a+\lambda$ changes $b'$ by
$-\theta B_0\lambda$.  There are finitely many faces
in the fundamental polyhedra.  Cut their images at the
integer phase hyperplanes and subdivide; this makes
every face map a finite piecewise affine map.  Primitive
edge-circle quotients are the integral linear maps
$b\mapsto\det(n_k,b)\bmod\mathbb Z$.
Triangulate each occurrence of a torus separately and
interpolate through a common subdivision along each
mapping cylinder.  This does not require a finite
triangulation invariant under the infinite monodromy
group.  Products of simplices with intervals are
triangulated by their ordered vertex prisms.  These
rules supply finite characteristic maps for \eqref{maps:eq:B.5.8},
and pulling the cells back gives the asserted CW
structure of $X$.

\item For the smaller homotopy model put the usual product
cell structure on each torus and apply cellular
approximation to the displayed maps.  This changes
the double mapping cylinder by homotopy equivalence;
equivalently use \cite[Propositions 4G.1--4G.2]{mapsHatcher02}.
The two finite surfaces are mapping tori of $T_L$,
because $m_j\gamma$ is a primitive circle coordinate
on $S_j$; their linear monodromies are
$B_j=A_j|_L$ (an inverse convention gives the same
counts).  Their cell counts are $(1,4,6,4,1)$.
The counts for $U_*$ are $(1,6,14,16,9,2)$, and
each $M_j$ has counts $(1,5,10,10,5,1)$.
In \eqref{maps:eq:B.5.4} the quotient honeycomb has two vertices,
three edges, and one hexagon.  The surviving phase
spaces are respectively a point, a circle, and a
two-torus, giving counts $(2,3,4,2,1)$ for $W$.
Adding these cells and one prism cell for each
boundary cell gives \eqref{maps:eq:B.5.2}.
\end{enumerate}
\end{proof}

\begin{theorem}\label{maps:src:B.5.2}
\StatementRecord{maps:src:B.5.2}{proved}{Historical events: Claude Fable 5.2 (Anthropic), effort high, 16 November 2026: carried as background without fault found. GPT-6 Astra (OpenAI), ChatGPT Pro mode, effort Pro, S6\_09\_PAPERS, 17 November 2026: targeted proof check. Also checked in this writing. Pre-continuation writing events: GPT-6 Astra (OpenAI), ChatGPT Pro mode, effort Pro, S6\_10\_PAPERS, 18 November 2026: targeted proof check. Continuation metadata events: GPT-6 Astra (OpenAI), ChatGPT Pro mode, effort Pro, S6\_10\_PAPERS continuation, 19 November 2026: used and rechecked in this run; status-vocabulary reconciliation; metadata reconciliation only; no new mathematical check. \emph{Pre-continuation recorded proof scope.} GPT-6 Astra (OpenAI, ChatGPT Pro mode), effort Pro, 2026-11-17; Claude Fable 5.2 (Anthropic), effort high, 16 November 2026; also checked in this writing by GPT-6 Astra (OpenAI), ChatGPT Pro mode, effort Pro, 18 November 2026}{the hypotheses, notation and definitions stated here; \cref{maps:src:B.5.1}; \cref{maps:eq:B.5.3}; \cref{maps:eq:B.5.6}; \cref{maps:eq:B.5.11}; \cref{maps:eq:B.5.4}; \cref{maps:eq:B.5.15}; \cref{maps:eq:B.5.12}; \cref{maps:eq:B.5.18}; \cref{maps:eq:B.5.2}; \cref{maps:eq:B.5.9}; \cite[Theorems 4.32 and 4.5]{mapsHatcher02}}


The attaching data of Theorem \ref{maps:src:B.5.1} give
$\pi_1(X)=1$ and integral homology
$H_k(X;\mathbb Z)=\mathbb Z$ for $k=0,6$ and zero
otherwise.  They give an explicit recognition of $X$
as a homotopy sphere, and exhibit the normally framed
three-torus used in B.1.

\end{theorem}


\begin{proof}
\begin{enumerate}[label=\arabic*.,leftmargin=*]
\item The graph piece has presentation
\[
 \pi_1(U_*)=
 \langle\Lambda,\rho_1,\rho_2\mid
 [\Lambda,\Lambda]=1,\quad
 \rho_jt_\lambda\rho_j^{-1}=t_{A_j\lambda}\rangle.
\]
Formula \eqref{maps:eq:B.5.3} imposes $\rho_j^{m_j}=t_{v_j}$.
Formula \eqref{maps:eq:B.5.6} kills $w,d$ and the cusp meridian
$a_0=\rho_2^{-1}\rho_1^{-1}$.  Since $A_1w=u-w$,
the normal closure also kills $u$.  Writing $c$ for
the image of $g$, the remaining presentation is
\[
 \langle c,x,y\mid c\text{ central},\quad
 xy=1,\quad x^3=c,\quad y^4=c^{-1}\rangle.
 \tag{B.5.9}\label{maps:eq:B.5.9}
\]
Substituting $y=x^{-1}$ gives $x^3=x^4$, so $x=c=y=1$.
This calculates the nonabelian fundamental group,
not only its abelianization.

\item We record the integral homology maps in primitive
bases.  Write $\gamma,u^*,w^*,d^*$ for the cohomology
basis dual to $g,u,w,d$, and put
\[
 \eta_1=2u^*+w^*+3d^*,\quad
 \eta_2=u^*+w^*+2d^*,\quad
 q_2=u^*w^*+6\gamma d^*.
 \tag{B.5.10}\label{maps:eq:B.5.10}
\]
Products here are exterior products.  First attach
the two finite pieces over a disc containing $0,1$;
call the result $X^\circ$.  Its homotopy type is the
double mapping cylinder $S_1\leftarrow F\to S_2$.
The degree-one
map $(\pi_{1*},-\pi_{2*})$ is
\[
 \alpha_1=
 \begin{pmatrix}
 3&0&0&0\\0&2&1&3\\4&0&0&0\\0&-1&-1&-2
 \end{pmatrix}.
 \tag{B.5.11}\label{maps:eq:B.5.11}
\]
To verify its primitive target bases, abelianize
$\pi_1(S_j)=\langle\Lambda,\rho_j\mid
\rho_j\lambda\rho_j^{-1}=A_j\lambda,
\rho_j^{m_j}=t_{v_j}\rangle$.
For $S_1$ the images are
$g=3\rho_1,u=2w,d=3w$; for $S_2$ they are
$g=-4\rho_2,u=w,d=2w$.  Thus
$(\rho_1,w)$ and $(\rho_2,w)$ are free integral bases.
This proves \eqref{maps:eq:B.5.11} without a rational invariant-lattice
identification.  Its Smith form is $(1,1,1,0)$,
its kernel is generated by
$\nu_0=-u-w+d$, and its primitive left kernel is
$(4,0,-3,0)$.  Consequently
$H_1(X^\circ)=\mathbb Z\zeta$, normalized by
\[
 g\mapsto12\zeta,\quad
 \rho_1\mapsto4\zeta,\quad
 \rho_2\mapsto-3\zeta,\quad
 a_0\mapsto-\zeta.
 \tag{B.5.12}\label{maps:eq:B.5.12}
\]

\item For completeness we derive the finite degree-two
pullback lattices.  The mapping-torus matrices on
$L$ are
\[
 B_1=\begin{pmatrix}0&1&0\\-1&-1&0\\1&0&1\end{pmatrix},
 \qquad
 B_2=\begin{pmatrix}0&-1&0\\1&0&0\\0&1&1\end{pmatrix}.
\]
The exterior-power mapping-cone complexes show that
$H_*(S_j)$ is free, with ranks $(1,2,2,2,1)$:
the nonzero Smith entries of
$\bigwedge^r B_j-I$ are all one.  In particular the
intersection lattice in degree two is unimodular.
Direct exterior multiplication gives the primitive
invariant lattices
$I_1=\langle\gamma\eta_1,q_2\rangle$ and
$I_2=\langle\gamma\eta_2,q_2\rangle$.
Their intersection determinants have absolute values
$9$ and $4$.  The pullback under a covering of degree
$m_j$ multiplies the intersection pairing by $m_j$.
Its image therefore has determinant $m_j^2$ and index
one in $I_1$, respectively index two in $I_2$.
The embedded torus generated by $v_2,d$ is invariant
under $h_2$ and covers its quotient with degree four.
The class $a\gamma\eta_2+bq_2$ evaluates on it as
$-2a-6b$.  A pullback class must therefore satisfy
$a+b\equiv0\pmod2$.  This is already an index-two
condition and hence determines the image exactly.
Thus
\[
 \pi_1^*H^2(S_1)=\langle\gamma\eta_1,q_2\rangle,
 \qquad
 \pi_2^*H^2(S_2)=\langle2\gamma\eta_2,
                                  \gamma\eta_2+q_2\rangle.
 \tag{B.5.13}\label{maps:eq:B.5.13}
\]

\item Use bases with pullbacks
$(\gamma\eta_1,q_2-2\gamma\eta_1)$ and
$(2\gamma\eta_2,q_2-3\gamma\eta_2)$, and their
Kronecker-dual homology bases.  Evaluating on
$(gu,gw,gd,uw,ud,wd)$ gives
\[
 \alpha_2=
 \begin{pmatrix}
 2&1&3&0&0&0\\-4&-2&0&1&0&0\\
 -2&-2&-4&0&0&0\\3&3&0&-1&0&0
 \end{pmatrix}.
 \tag{B.5.14}\label{maps:eq:B.5.14}
\]
Its Smith form is again $(1,1,1,0)$ and its primitive
left kernel is $\Phi=(4,2,3,2)$.  Thus its cokernel is
$\mathbb Z e$, where the cohomology class $\mathcal A$
given by $\Phi$ satisfies
$\langle\mathcal A,e\rangle=1$ and restricts on $F$
to $2q_2$.  The exact sequence of this double mapping
cylinder gives
\[
 0\longrightarrow\mathbb Z e\longrightarrow
 H_2(X^\circ)\longrightarrow\mathbb Z\nu_0
 \longrightarrow0.
 \tag{B.5.15}\label{maps:eq:B.5.15}
\]

\item The central quotient \eqref{maps:eq:B.5.4} has cellular boundary
$\partial_1(e_k)=Q-P$ and all higher cellular
boundaries zero.  Indeed the hexagon attaching word
can be chosen as
$e_1e_2^{-1}e_3e_1^{-1}e_2e_3^{-1}$; every edge occurs
with both signs.  Crossing with a surviving phase
circle preserves this cancellation, and the torus
cell boundary is the commutator.  Hence
\[
 H_*(W)=(\mathbb Z,\mathbb Z^2,\mathbb Z^4,
                       \mathbb Z^2,\mathbb Z).
\]
At $\theta=0$ the cusp map is the quotient
$c(a,b)=[a,b]$.  On $H_1$ it is projection to
$\langle g,u\rangle$.  With degree-two target basis
$(T_B,\overline C_1,\overline C_2,\overline C_3)$,
its matrix is
\[
 c_2=\begin{pmatrix}
 1&0&0&0&0&0\\0&0&-1&0&0&0\\
 0&-1&-1&1&1&0\\0&0&0&1&0&0
 \end{pmatrix}.
 \tag{B.5.16}\label{maps:eq:B.5.16}
\]
For a mixed torus $\beta v$ this follows from the
degree formula
$\sum_k\det(n_k,B_0\beta)\det(n_k,v)[\overline C_k]$.
The first factor is its oriented honeycomb edge
coefficient, and the second is the degree of the
primitive circle quotient.  The map sends $gu$ to
$T_B$ and $wd$ to zero.  In degree three it sends
$guw,gud$ to a basis and kills $gwd,uwd$; in degree
four it has degree one on the open hexagon times
$T_f$.  All these maps are integrally surjective.
Their kernels are exactly
$\operatorname{im}(\bigwedge^r M_0-I)$: inclusion
follows by transport around the boundary, and the
displayed kernels and monodromy images have the
same primitive bases, as checked in
\texttt{code/maps\_cw.py}.

\item Let $K_i$ be the kernel of $H_i(M_0)\to H_i(W)$.
The mapping-torus sequence and the preceding
surjections give
\[
 K_1=\mathbb Z a_0,\qquad
 K_2=\mathbb Z\xi_w\oplus\mathbb Z\xi_d,
 \qquad \xi_v=[S^1_v\times S^1_{a_0}].
 \tag{B.5.17}\label{maps:eq:B.5.17}
\]
These are literal tori because $M_0$ fixes $w,d$.
Under the mapping-torus projection they go to $w,d$,
so their span is primitive.  In \eqref{maps:eq:B.5.6}, over a
hexagon centre the $a_0$ coordinate is collapsed;
thus their images in $H_2(W)$ vanish.  The ranks
from the mapping-torus sequence show that this is
the whole kernel.

\item Orient the connecting map in \eqref{maps:eq:B.5.15} so that
\[
 \partial j_*\xi_v=(I-A_2)v
                   =(I-A_1^{-1})v\quad(v\in\langle w,d\rangle).
\]
One verifies this on the two rectangles obtained
by cutting the swept torus along the finite band:
their unmatched boundary circles are $v$ and
$-A_2v$.  Since $(I-A_2)w=-\nu_0$ and
$(I-A_2)d=0$, the class $z=j_*\xi_w$ together with
$e$ is a basis of $H_2(X^\circ)$.
To determine $j_*\xi_d$, use the circle translation
in direction $d$ on $X^\circ$.  For this action the
integral sweep operator on cohomology restricts on
$F$ to contraction by $d$.  The class
$\mathcal A|_F=2q_2$ satisfies
\[
 \iota_d(2q_2)=-12\gamma.
\]
Restriction $H^1(X^\circ)\to H^1(F)$ is injective
by \eqref{maps:eq:B.5.12}, with image $\mathbb Z\,12\gamma$.
Thus $\iota_d\mathcal A=-\Theta$, where
$\Theta(\zeta)=1$.  Evaluating the swept torus and
using $\Theta(a_0)=-1=p$ gives
\[
 j_*\xi_d=-p e=e,
 \qquad
 [j_*|_{K_2}]_{(\xi_d,\xi_w)}^{(e,z)}
       =\begin{pmatrix}1&0\\0&1\end{pmatrix}.
 \tag{B.5.18}\label{maps:eq:B.5.18}
\]
This computes the attaching coefficient that cannot
be recovered from ranks alone.

\item In degree one the full cusp attachment map in bases
$(g,u,a_0)$ and $(\bar g,\bar u,\zeta)$ is
\[
 D_1=\begin{pmatrix}1&0&0\\0&1&0\\-12&0&1\end{pmatrix}.
\]
In degree two, split the surjection onto $H_2(W)$;
the remaining block is exactly \eqref{maps:eq:B.5.18}.  Thus the
full map is integrally equivalent to $I_4\oplus I_2$.
Both are unimodular.  Mayer--Vietoris proves
$H_1(X)=H_2(X)=0$, with no torsion.  The cell count
\eqref{maps:eq:B.5.2} gives $\chi(X)=2$.  Since $X$ is a closed
oriented six-manifold, duality and universal
coefficients give $H_4=H_5=0$ and make $H_3$ free:
$\operatorname{Ext}(H_3,\mathbb Z)$ injects into
$H^4(X)=H_2(X)=0$.  The Euler characteristic then
gives $\operatorname{rank}H_3=0$.  This proves all
the asserted integral groups.  The program checks
every matrix, Smith form, primitive coefficient
and cell count just used, using exact arithmetic.

\item Starting from \eqref{maps:eq:B.5.9}, successive applications of
the Hurewicz theorem make $\pi_i(X)=0$ for
$1\leq i\leq5$.  The same theorem produces a map
$S^6\to X$ carrying the fundamental class to the
chosen generator.  The mapping-cylinder pair of
this map is simply connected and has zero relative
homology.  If it had a first nonzero relative
homotopy group, the relative Hurewicz theorem
would make its relative homology nonzero.  All
relative homotopy groups therefore vanish, and
the CW Whitehead theorem makes the map a homotopy
equivalence.  These
are \cite[Theorems 4.32 and 4.5]{mapsHatcher02}
applied to the explicitly computed complex.

\item Finally keep a small regular disc in the pair of
pants uncollapsed while forming its ribbon spine.
The same construction relative to that disc retains
the block $F\times D^2$ with its flat vertical
frame and base frame.  The three-torus
$\{\gamma=\theta_0\}\subset F$ has normal frame
given by the two base directions and the
$\gamma$ direction.  This is exactly the regular
fibre and normal framing of B.1.  The CW
presentation therefore retains the geometric
framed fibre; no smooth framing of hypothetical
cancelled handles has been inferred from the
homology calculation.
\end{enumerate}
\end{proof}

